\subsection{Apéndices}
\label{subsec:apendices}
% ####################################################################################################################

\subsubsection{Significado de los parámetros del ajuste lineal ($m$ y $b$)}
\label{subsubsec:significado-parametros-ajuste-lineal}
Al confrontar el modelo teórico de la Ley de Ohm ($V = I \cdot R$) con la ecuación geométrica de la recta ajustada ($V = m \cdot I + b$), se realiza una identificación unívoca de sus componentes. La pendiente $m$ se corresponde directamente con el valor de la resistencia equivalente del conductor bajo ensayo ($R = m$). Por su parte, la ordenada al origen $b$ representa el valor extrapolado de la diferencia de potencial cuando la corriente del sistema es nula. 

Físicamente, aunque en un escenario ideal se espera que $b = 0$, en el laboratorio real suelen presentarse ligeras desviaciones ($b \neq 0$). Estas discrepancias sistemáticas no reflejan una falla del modelo óhmico, sino que se atribuyen de forma directa a factores instrumentales y de montaje. Entre ellos destaca una pequeña descalibración interna de la fuente de alimentación que hace que entregue un voltaje mínimo residual incluso cuando su perilla está configurada exactamente en cero (fenómeno técnicamente denominado tensión de \textit{offset}), así como la existencia de resistencias de contacto en las superficies de fricción donde las pinzas cocodrilo muerden los bornes de la resistencia cerámica.

En \textbf{conclusión}, la pendiente $m$ equivale a la resistencia del conductor ($R$) y la ordenada $b$ representa el potencial residual a corriente cero, cuyo desvío respecto a un valor nulo se debe a imperfecciones de los instrumentos y a resistencias de contacto en los bornes.

\subsubsection{Análisis comparativo: Regresión lineal global frente a promedio punto a punto}\label{subsubsec:comparacion-regresion-lineal-vs-promedio-punto-a-punto}
La propuesta de determinar la resistencia mediante el cálculo individual de los cocientes $R_n = V_n / I_n$ para cada par de datos medidos y luego calcular su promedio aritmético simple no es metodológicamente equivalente al ajuste por cuadrados mínimos. El cálculo individual propaga y amplifica de manera directa las incertidumbres absolutas asociadas a las lecturas de los instrumentos en cada punto de ensayo. 

Este fenómeno es severamente crítico en el régimen de bajas corrientes (las primeras mediciones del ensayo), donde los valores de intensidad se aproximan a la resolución nominal del amperímetro. En esta zona, pequeñas fluctuaciones aleatorias o el error de dígito del display generan variaciones desproporcionadas en el cociente $V_n / I_n$, dando lugar a valores artificialmente elevados o dispersos que sesgan y distorsionan fuertemente el promedio final. 

En contraposición, el método de cuadrados mínimos no promedia cocientes aislados, sino que evalúa la tendencia histórica y global de todo el conjunto de datos de forma simultánea. Al minimizar la suma de los residuos cuadráticos, la regresión pondera el comportamiento general del circuito, atenuando estadísticamente el impacto del ruido y de las incertidumbres de lectura en los regímenes más desfavorables del instrumento.


En \textbf{conclusión}, el promedio punto a punto deforma el resultado porque amplifica drásticamente los errores instrumentales en las lecturas de baja corriente, mientras que la regresión lineal por cuadrados mínimos provee una estimación estadísticamente robusta al mitigar las fluctuaciones individuales mediante la evaluación del conjunto total de datos.

% ####################################################################################################################
% ####################################################################################################################

\clearpage
\subsubsection{Cálculo detallado de incertezas de resistencias}
\label{subsubsec:calculo-incertezas-resistencias}

Se aplica la Ecuación~\ref{eq:incerteza-multimetro}: $\Delta R = \frac{\%rdg}{100} \times R_{medida} + dgt \times resolución$.

\textbf{R1} (rango 20 M$\Omega$: $\pm$2{,}0\% rdg, $\pm$2 dgt, resolución 10 k$\Omega$):

\begin{equation}
    \Delta R = 0{,}020 \times 5{,}61 + 2 \times 0{,}01 = 0{,}1122 + 0{,}02 = 0{,}1322 \ \text{M}\Omega \approx 0{,}1 \ \text{M}\Omega
    \label{eq:incerteza-R1}
\end{equation}

\begin{equation}
    \boxed{R_1 = (5{,}6 \pm 0{,}1) \ \text{M}\Omega}
\end{equation}

\textbf{R2} (rango 2000 k$\Omega$: $\pm$1{,}0\% rdg, $\pm$2 dgt, resolución 1 k$\Omega$):

\begin{equation}
    \Delta R = 0{,}010 \times 562 + 2 \times 1 = 5{,}62 + 2 = 7{,}62 \ \text{k}\Omega \approx 8 \ \text{k}\Omega
    \label{eq:incerteza-R2}
\end{equation}

\begin{equation}
    \boxed{R_2 = (562 \pm 8) \ \text{k}\Omega}
\end{equation}

\textbf{R3} (rango 20 k$\Omega$: $\pm$0{,}7\% rdg, $\pm$1 dgt, resolución 10 $\Omega$):

\begin{equation}
    \Delta R = 0{,}007 \times 5500 + 1 \times 10 = 38{,}5 + 10 = 48{,}5 \ \Omega \approx 50 \ \Omega = 0{,}05 \ \text{k}\Omega
    \label{eq:incerteza-R3}
\end{equation}

\begin{equation}
    \boxed{R_3 = (5{,}50 \pm 0{,}05) \ \text{k}\Omega}
\end{equation}

% ####################################################################################################################
% ####################################################################################################################

\clearpage
\subsubsection{Modelo del divisor resistivo con voltímetro real}
\label{subsubsec:divisor-resistivo-voltimetro-real}

\todo{Deducción de $V_{le\acute{i}da}$ con $R'_2 = R_2 \| R_{int}$; cálculo numérico para $R_1=R_2 \in \{5{,}6\ \text{k}\Omega;\ 560\ \text{k}\Omega;\ 5{,}6\ \text{M}\Omega\}$ comparado con el valor ideal.}

% ####################################################################################################################
% ####################################################################################################################

\clearpage
\subsubsection{Preguntas de la guía experimental}
\label{subsubsec:preguntas-guia}

\todo{Responder preguntas}

\paragraph{Práctica 1 --- Ley de Ohm}
\begin{itemize}
    \item ¿Qué significado tienen los valores de la pendiente $m$ y la ordenada al origen $b$?
    
    \respuesta{La pendiente $m$ equivale física y unívocamente a la resistencia eléctrica del conductor metálico ($R = m$), mientras que la ordenada al origen $b$ representa el potencial eléctrico residual en vacío cuando la corriente eléctrica es nula ($I = 0$). El desglose físico de estos parámetros y las causas de la desviación experimental de la ordenada respecto a cero ($b \neq 0$) se detallan en la sección \ref{subsubsec:significado-parametros-ajuste-lineal}.}

\item ¿Qué tal si para cada par $(V_n, I_n)$ calculamos $R_n = V_n/I_n$ y luego promediamos los $R_n$? ¿Estamos haciendo lo mismo que la regresión? ¿Qué daría en nuestro caso?
    
    \respuesta{No es metodológicamente equivalente. El cálculo individual punto a punto amplifica drásticamente las incertidumbres instrumentales absolutas, afectando de manera crítica el régimen de bajas corrientes debido a la resolución nominal del amperímetro, lo cual deforma el promedio final. Por el contrario, la regresión lineal evalúa de manera robusta la totalidad de los datos mitigando estadísticamente las fluctuaciones aleatorias individuales. La justificación y el análisis físico detallado de esta diferencia fundamental se presentan en la sección \ref{subsubsec:comparacion-regresion-lineal-vs-promedio-punto-a-punto}.}
\end{itemize}

\paragraph{Práctica 3 --- Influencia del aparato de medida}
\begin{itemize}
    \item ¿Qué resultado se espera al sumar (con signo) las caídas de tensión en el circuito? ¿Qué se obtuvo?

    \respuesta{Completar.}

    \item ¿Qué sucede con la Segunda Ley de Kirchhoff? ¿A veces anda y a veces no?

    \respuesta{Completar.}
\end{itemize}